% =========================
% kolory
% =========================
\definecolor{TaskGray}{HTML}{F7F7F8}

% =========================
% globalna zwartość / odstępy
% =========================
\setlength{\parindent}{0pt}
\setlength{\parskip}{1.6pt}
\setlength{\columnsep}{0.9cm}
\setlength{\columnseprule}{0.5pt}
\linespread{1.02}

% =========================
% itemize, steps i podpunkty 
% =========================

\setlist[itemize]{
  label=\textbullet,
  leftmargin=1.2em,
  itemsep=0pt,
  topsep=0pt,
  parsep=0pt,
  partopsep=0pt
}

\newlist{steps}{enumerate}{1}
\setlist[steps]{
  label=\textbf{(\arabic*)},
  leftmargin=1.9em,
  itemsep=0.42em,
  topsep=0.2em,
  parsep=0pt,
  partopsep=0pt
}

\newlist{podpunkty}{enumerate}{1}
\setlist[podpunkty]{
  label=\textbf{(\alph*)},
  leftmargin=1.9em,
  itemsep=0.42em,
  topsep=0.2em,
  parsep=0pt,
  partopsep=0pt
}


% =========================
% indeks
% =========================
\newcommand{\indeks}[1]{%
  \par\bigskip
  \noindent
  {\setlength{\fboxsep}{8pt}%
   \setlength{\fboxrule}{1.3pt}%
   \fcolorbox{black}{white}{%
     \strut\fontsize{22}{26}\selectfont\bfseries #1%
   }}%
  \par\vspace{0.95em}
}

% =========================
% temat i bit
% =========================
\newcommand{\temat}[1]{%
  \par\medskip
  \noindent{\bfseries
      \fontsize{16}{20}\selectfont
      #1%
  }
  \par\medskip
}

\newcommand{\bit}[1]{%
  \par\medskip
  \noindent{\bfseries #1}%
  \par\smallskip
}

% =========================
% logbox
% =========================
\newtcolorbox{logbox}{
  enhanced,
  colback=TaskGray,
  colframe=white,
  boxrule=0.9pt,
  arc=0pt,
  left=6pt,
  right=6pt,
  top=5pt,
  bottom=5pt,
  boxsep=0pt,
  before skip=0.2em,
  after skip=0.7em
}

% =========================
% tw i def
% =========================
\newcommand{\twierdzenie}[2]{%
  \noindent
  #1
  \par\smallskip
  \noindent\hspace*{1.2em}$\Longrightarrow$
  \par\smallskip
  \noindent
  #2
  \par\smallskip
}

\newcommand{\definicja}[2]{\textbf{def #1}$\ \coloneqq\ $#2}

% =========================
% zadania (zwykłe i domowe)
% =========================
\newcommand{\zad}[1]{\textbf{z#1}}

\newcommand{\domek}[2]{%
  \tikz[baseline=(t.base)]{%
    \node[inner xsep=2.5pt, inner ysep=1.8pt, font=\bfseries] (t) {#2};%
    \draw[line width=0.6pt, line join=round, #1]
      (t.south west) -- (t.north west) --
      ([xshift=-1.5pt]t.north west) -- ([yshift=1.1ex]t.north) --
      ([xshift=1.5pt]t.north east) -- (t.north east) --
      (t.south east) -- cycle;%
  }%
}

\newcommand{\zd}[1]{\domek{}{z#1}}                              % domowe z numerem (pod kolokwium)
\newcommand{\zdc}{\domek{dash pattern=on 1.5pt off 1pt}{z}}    % domowe bez numeru (ciekawostka)
\newcommand{\itemzd}{\stepcounter{podpunktyi}\item[\domek{}{(\alph{podpunktyi})}]}   % domowy podpunkt w podpunkty